% Preámbulo común de los informes E1, E2 y E3, con aspecto de Word: Calibri,
% paleta del tema de Office, bandas y portada.
% Cada informe lo carga con \input{../../latex/preambulo.tex}.
%   \portada[Alumnos]{sigla y curso}{tipo de entrega}{título}{autores}{fecha}
\usepackage[spanish,es-tabla]{babel}
\usepackage[utf8]{inputenc}
\usepackage[T1]{fontenc}
% Carlito: mismas medidas que Calibri, la fuente de Word
\usepackage[sfdefault,lining]{carlito}
\usepackage[letterpaper,margin=2.2cm,top=2.8cm,bottom=2.8cm]{geometry}
\usepackage{xcolor}
\usepackage{graphicx}
\usepackage{array}
\usepackage{colortbl}
\usepackage{tabularx}
\usepackage{microtype}  % reparte el justificado en letras y espacios: menos huecos en columnas angostas
\usepackage{ragged2e}
\setlength{\JustifyingParindent}{0pt}  % sin sangría, como el resto del texto
\usepackage{enumitem}
\usepackage{tikz}
\usetikzlibrary{positioning,arrows.meta,calc,shapes}
\usepackage{eso-pic}
\usepackage{titlesec}
\usepackage{listings}
\usepackage[hidelinks]{hyperref}
\graphicspath{{../../assets/}}

% paleta del tema de Office
\definecolor{banda}{HTML}{4F81BD}
\definecolor{titulo}{HTML}{2E75B5}
\definecolor{subtitulo}{HTML}{5B9BD5}
\definecolor{pizarra}{HTML}{273540}
\definecolor{tablaencabezado}{HTML}{DEEAF6}  % el azul claro de los diagramas
\definecolor{fondosuave}{HTML}{DEEAF6}
\definecolor{grisoffice}{HTML}{7F7F7F}
\definecolor{vinotinto}{HTML}{7B1E2B}

% texto agregado o cambiado por el asistente, pendiente de revisión del equipo
\newcommand{\nuevo}[1]{\textcolor{vinotinto}{#1}}

% títulos como los estilos de título de Word
\titleformat{\section}{\large\bfseries\color{subtitulo}}{}{0pt}{}
\titleformat{\subsection}{\normalsize\bfseries\color{subtitulo}}{}{0pt}{}
\titlespacing*{\section}{0pt}{2ex}{1ex}
\titlespacing*{\subsection}{0pt}{1.5ex}{0.6ex}
\setcounter{secnumdepth}{0}

\setlist{itemsep=0.4ex,leftmargin=1.6em}
\setlist[itemize]{label=\textbullet}
\setlength{\parindent}{0pt}
\setlength{\parskip}{8pt plus 2pt}  % espacio después de párrafo de Word

% tablas tipo Word: grilla completa y encabezado con fondo
% celdas justificadas como el resto del texto; Y es la columna flexible de tabularx
\newcolumntype{L}[1]{>{\justifying\arraybackslash}p{#1}}
\newcolumntype{Y}{>{\justifying\arraybackslash}X}
\newcommand{\enc}[1]{\cellcolor{tablaencabezado}#1}
\newcommand{\iniciartabla}{\renewcommand{\arraystretch}{1.3}\arrayrulecolor{black!70}}

% notas al pie con números también dentro de bloques (minipage), como en Word
\renewcommand{\thempfootnote}{\arabic{mpfootnote}}

% pregunta que abre una respuesta
\newcommand{\pregunta}[1]{\textbf{\color{subtitulo}#1}}

% SQL y pseudocódigo
\lstset{basicstyle=\ttfamily\small,columns=fullflexible,keepspaces=true,
  frame=single,rulecolor=\color{grisoffice},backgroundcolor=\color{fondosuave},
  keywordstyle=\color{titulo}\bfseries,language=SQL,
  morekeywords={BEGIN,COMMIT,ROLLBACK},literate={á}{{\'a}}1 {é}{{\'e}}1
  {í}{{\'i}}1 {ó}{{\'o}}1 {ú}{{\'u}}1 {ñ}{{\~n}}1}

% portada: bandas completas; páginas interiores: media banda arriba a la
% izquierda, media banda abajo a la derecha con el número y logo de marca de agua
\pagestyle{empty}
\newif\ifenportada
\AddToShipoutPictureBG{%
  \ifenportada
    \AtPageUpperLeft{\color{banda}\raisebox{-1.2cm}{\rule{\paperwidth}{0.6cm}}}%
    \AtPageLowerLeft{\color{banda}\raisebox{0.6cm}{\rule{\paperwidth}{0.6cm}}}%
  \else
    \AtPageUpperLeft{\color{banda}\raisebox{-1.2cm}{\rule{0.62\paperwidth}{0.6cm}}}%
    \AtPageLowerLeft{%
      \put(\LenToUnit{0.45\paperwidth},\LenToUnit{0.6cm}){\color{banda}\rule{0.55\paperwidth}{0.6cm}}%
      \put(\LenToUnit{0.45\paperwidth},\LenToUnit{0.78cm}){\hspace{0.2cm}\color{white}\small\thepage}%
      \put(\LenToUnit{0.7cm},\LenToUnit{0.5cm}){\logomarcaagua}}%
  \fi}
\newcommand{\logomarcaagua}{\tikz\node[opacity=0.08,inner sep=0]{\includegraphics[width=2cm]{logo_uc.png}};}

% \portada[rótulo]{sigla y curso}{tipo de entrega}{título}{autores}{fecha}
\newcommand{\portada}[6][Alumno]{%
  \global\enportadatrue
  \begin{titlepage}
    \centering
    \vspace*{2.2cm}
    \includegraphics[width=8cm]{logo_uc.png}\par\vspace{1.6cm}
    {\large\bfseries\color{titulo} Magíster en Inteligencia Artificial - MIA\par}\vspace{0.4cm}
    {\large\bfseries\color{subtitulo} Instituto de Ingeniería Matemática y Computacional - IMC\par}\vspace{0.5cm}
    {\LARGE\bfseries\color{pizarra} #2\par}\vspace{0.5cm}
    {\itshape\bfseries\color{subtitulo} #3\par}\vspace{0.2cm}
    {\large #4\par}\vspace{0.2cm}
    {\color{banda}\rule{0.8\textwidth}{0.6pt}\par}
    \vspace{2cm}
    \raggedright
    {\textbf{#1:} #5\par}\vspace{0.3cm}
    {\textbf{Fecha de entrega:} #6\par}
    \vfill
  \end{titlepage}
  \global\enportadafalse}
